\section{Minimum Action as Survivorship}
\label{sec:minimum_action}
The Entropic Action defines \textit{what} a system must minimize; the Selection Principle establishes \textit{that} only minimal-action configurations remain observable. We now derive \textit{which} path a surviving system follows.

\subsection{The Exponential Filter}
\label{subsec:exponential_filter}
If persistence is maintained by state transitions, and each transition risks scrambling information into the substrate noise floor, then a history of $N$ transitions survives only if every step survives. Independent errors compound multiplicatively, yielding an exponential filter.

Consider a structural history $\Gamma$ accumulating entropic action $S_\Phi[\Gamma]$ (\cref{eq:entropic_action}). Because the substrate possesses a finite noise floor, each unit of entropic action carries
an intrinsic risk of information loss. Let $q$ be the probability that a single unit maintains causal coherence.

In QSS, successive units are statistically independent dissipative events. The probability that $\Gamma$ maintains coherence is therefore multiplicative:

\begin{equation}
\label{eq:survival_prob1}
    P(\text{survival} \mid \Gamma) 
    = q^{S_\Phi[\Gamma]} 
\end{equation}

Using the identity $q^C = e^{C\ln q}$ with $C = S_\Phi[\Gamma]$, and defining $\beta \equiv -\ln q > 0$ (positive since $0 < q < 1$), we have $S_\Phi[\Gamma]\ln q = -\beta S_\Phi[\Gamma]$, turning \cref{eq:survival_prob1} into:

\begin{equation}
\label{eq:survival_prob}
    P(\text{survival} \mid \Gamma) = e^{-\beta S_\Phi[\Gamma]}
\end{equation}

Here $\beta \equiv -\ln q$ is the substrate's inverse entropic tolerance, a geometric invariant of the channel. For discrete substrates, $q$ is determined by the channel geometry; in continuous media, $\beta$ emerges from the substrate's noise spectral density.

The form itself is the informational analogue of the Boltzmann distribution $P \propto e^{-E/k_BT}$, with $S_\Phi$ playing the role of energy and $\beta$ the role of inverse temperature. The Boltzmann form is derived from ensemble statistics at equilibrium and this form is derived from the multiplicative compounding of independent transition risk along a single history, presupposing neither equilibrium, a thermal bath, nor a physical substrate. Both converge on the same exponential law; the informational route recovers it under strictly weaker assumptions: the substrate's own noise geometry does the work that equilibrium, a thermal bath, and a variational functional do for the classical derivation.

\subsection{The Severity of the Selection Principle}
\label{subsec:severity}

How harsh is the Selection Principle? Let $\Gamma_{\text{min}}$ be the minimal-action trajectory and $\Gamma_{\text{var}}$ any variant path with a deviation $\delta S_\Phi > 0$. Over any observation window, the relative survival probability is:

\begin{equation}
\label{eq:severity}
    \frac{P(\Gamma_{\text{var}})}{P(\Gamma_{\text{min}})} = e^{-\beta \delta S_\Phi}
\end{equation}

When the fundamental update cycle is much shorter than the observation timescale, even a small per-transition deviation compounds into a vast accumulated $\delta S_\Phi$, and the survival ratio $P(\Gamma_{\text{var}})/P(\Gamma_{\text{min}})$ falls to zero. The gap between the best trajectory and the next-best widens without bound.

\subsection{No Merit Required}
\label{subsec:teleology}
Severity is a statement about the losers, not the winner. It says nothing about whether $\Gamma_{\text{min}}$ is efficient, elegant, or well-suited to anything, only that every alternative to it decayed faster. A mediocre trajectory survives just as certainly as an excellent one, provided everything else was worse. What looks like design, foresight, or purpose is not a property of the winner; it is the absence of its competition.

\subsection{Minimum Action as Survivorship}
\label{subsec:minimum_action_as_survivorship}
The severity of the Selection Principle is formalized as:

\begin{proposition}[Minimum Action as Survivorship]
In QSS, the observable trajectories of any persistent system are those achieving the minimal Entropic Action $S_\Phi$ among all candidate trajectories.
\end{proposition}

\begin{proof}
By the Selection Principle and \cref{eq:severity}, only trajectories minimizing $S_\Phi$ remain observable; the observable set is therefore the argmin of $S_\Phi[\Gamma]$ over candidate histories $\Gamma$, possibly containing more than one history if the minimum is not unique.
\end{proof}

Standard quantum mechanics grounds the Principle of Least Action in phase interference. The Selection Principle offers a distinct, classical mechanism (informational elimination) by which nature could favor certain histories: the substrate itself acts as a filter, and what remains observable is only what has not yet dissipated.